\section{Oppdragsgivar}
\textbf{Renasys AS i samarbeid med Sunnfjord kommune}

Renays \citep{Renasys} er ei innovativ og nyskapande oppstartsbedrift som arbeider med banebrytande teknologi innan mekanisk finpartikkelfiltrering av avlaupsvatn.
Bedrifta har 15 tilsette fordelt på forskjellige kontor. Dei har kontor på Øyrane i Førde og Sandnes i Rogaland. 
Etter å ha arbeidd konfidensielt over lengre tid gjekk Renasys offentleg ut med teknologien sin i løpet av 2023. 
Dei tilbyr reinsetenester til kommunar og interkommunale selskap innan avlaup og maritime sektor.

Samarbeidet med \gls{Sunnfjord Kommune} \citep{SunnfjordKommune} er retta mot ``Mission Zero'' som er eit ambisiøst mål om 
null utslepp, null avfall, null energi og ein generell forbetring av avlaupssektoren i Noreg. \newline
Sunnfjord kommune er ansvarleg for vann, veg og avlaup i sitt område og har engasjert 
Renasys for å utforske forbetringar ved reinseanlegget på Sande.
\newline
\newline


\begin{figure}[htbp]
    \centering
    \begin{subfigure}[b]{0.3\textwidth}
        \centering
        \includegraphics[width=1\textwidth]{Bilder/renasys.png}
    \end{subfigure}
    \hfill
    \begin{subfigure}[b]{0.3\textwidth}
        \centering
        \includegraphics[width=0.7\textwidth]{Bilder/SK.png}
    \end{subfigure}
    \caption{Logo oppdragsgivar}\label{fig:Oppdragsgivar}
\end{figure}

\section{Vedlegg}
Til denne rapporten er det utarbeidd eit eige dokument som inneheld alle relevante vedlegg. 
Desse vedlegga omfattar ytterlegare tekniske detaljar knytt til oppgåva vår.\newline
Vedlegga er organiserte for å gje lesaren ein djupare innsikt i arbeidet vi har utført.